\chapter{Outlook}
\label{chap:outlook}

The sensitivity comparison reported in this thesis is one stage of an analysis that aims to publish constraints on the dimension-eight aQGC Wilson coefficients in the fully hadronic VBS channel. Three steps remain: finalising the signal region (Section~\ref{sec:outlook_signalregion}), setting limits on the Wilson coefficients and placing them in context with the existing measurements (Section~\ref{sec:outlook_limits}), and extending the measurement to the Run-3 dataset (Section~\ref{sec:outlook_run3}).

\section{Finalisation of the signal region}
\label{sec:outlook_signalregion}

The training selection used throughout this thesis (Section~\ref{sec:event_selection}) preserves enough background statistics for the MLP to converge while bringing the phase space close to the analysis signal region. A publishable analysis requires a tighter signal region that enables a data-driven background estimate and maps cleanly onto the $WW$, $WZ$ and $ZZ$ decay channels. Two refinements remain.

First, the DisCoJet${}>0.65$ requirement on both large-$R$ signal jets, shown in Section~\ref{sec:event_selection} to outperform the standard tight ATLAS working point of $0.85$ on the dominant tensor operators, must be re-evaluated under the constraint that the residual QCD multijet contamination stays compatible with a data-driven sideband estimate; splitting the signal region into orthogonal DisCoJet categories may additionally expose the boson-decay purity as a handle for the profile-likelihood fit. Second, the large-$R$ jet mass, currently only an MLP input feature, requires an explicit window around the $W$ and $Z$ pole masses. Such a window enables the sideband-based QCD multijet estimate that would replace the conservative $30\%$ envelope of Section~\ref{sec:disc_systematics}, and provides the natural definition of the $WW$, $WZ$ and $ZZ$ categories, since the $\approx 10\,\mathrm{GeV}$ pole separation is resolvable in the large-$R$ jet mass. Because DisCoJet is decorrelated from the jet mass, the window can be imposed without sculpting the background. Both choices are fixed by the same per-operator Asimov-significance scan used in Section~\ref{sec:event_selection}.

\section{Limit setting on the Wilson coefficients}
\label{sec:outlook_limits}

The Asimov significances reported in this thesis are projected onto $95\%$ confidence-level limits on the Wilson coefficients through a profile-likelihood fit to the binned MLP discriminant template. Because the signal yield scales as $(f_i/\Lambda^4)^2$ in the quadratic-only model, the limit on a coefficient improves with integrated luminosity, fastest for the signal-limited tensor operators and more slowly for the background-limited scalar and mixed operators (Sections~\ref{sec:sens_scalar} and~\ref{sec:sens_tensor}).

The strongest existing constraints on the Éboli basis come from Run-2 VBS measurements at $\sqrt{s}=13$~TeV. The individual inputs are the ATLAS~\cite{ATLAS:semileptonicVBS:2025} and CMS~\cite{CMS:semileptonicVBS:2022} semi-leptonic $WV$ searches and the fully-leptonic $ZZjj$ analyses of ATLAS~\cite{ATLAS:fullyleptonicVBS:2024} and CMS~\cite{CMS:fullyleptonicZZjj:2021}. The tightest published limits are obtained from the ATLAS combined EFT interpretation of seven vector-boson-scattering measurements and one triboson measurement~\cite{ATLAS:VBStribosoncombination:2026}, which exploits the correlations between channels. Its observed, unitarity-bounded intervals for the operators probed in this thesis are collected in Table~\ref{tab:best_limits}. The constraints span more than two orders of magnitude across the families: the tensor coefficients are bounded at the $\mathcal{O}(0.1\text{--}0.6)\,\mathrm{TeV}^{-4}$ level, the mixed coefficients between $\mathcal{O}(2)$ and $\mathcal{O}(20)\,\mathrm{TeV}^{-4}$, and the scalar sector only at the $\mathcal{O}(5)\,\mathrm{TeV}^{-4}$ level, mirroring the cross-section hierarchy that makes the scalar family the least sensitive (Section~\ref{sec:sens_scalar}). The intervals quoted include unitarity bounds: because the dimension-eight amplitude grows as $E^4/\Lambda^4$, a sufficiently large coefficient would drive the predicted $VV\to VV$ amplitude past the limit set by partial-wave unitarity, where the EFT no longer describes a physical theory, so the amplitude is capped to keep the extracted limit meaningful. The bounded intervals are therefore more conservative than the corresponding unbounded ones, which are numerically tighter but rely on an unphysical high-energy extrapolation.

% Values from Table 2 of arXiv:2603.18630 (ATLAS VBS+triboson combination): observed 95% CL,
% one coefficient at a time, WITH unitarity bounds. S0 = combined S0+S2 direction (f_S02);
% S1 has no unitarised interval; T3, T4 not in the combination. Read from the arXiv HTML --
% verify against the published Table 2 / HEPData before final submission.
\begin{table}[htbp]
  \centering
  \caption{Best published $95\%$~CL constraints on the dimension-eight Wilson coefficients $f_i/\Lambda^4$ relevant to this analysis, in $\mathrm{TeV}^{-4}$, from the ATLAS combined EFT interpretation of vector-boson scattering and triboson measurements at $\sqrt{s}=13$~TeV with $140~\mathrm{fb}^{-1}$~\cite{ATLAS:VBStribosoncombination:2026}. Intervals are observed, obtained one coefficient at a time with unitarity bounds applied. $S_0$ denotes the combined $\mathcal{O}_{S,0}{+}\mathcal{O}_{S,2}$ direction varied with a single coefficient $f_{S0,2}$ (Chapter~\ref{chap:sm}); $S_1$ ($-$) has no unitarised interval, and $T_3$, $T_4$ are not included in the combination.}
  \label{tab:best_limits}
  \begin{tabular}{|cc|cc|cc|}
    \hline
    \multicolumn{2}{|c|}{\textbf{Scalar}} & \multicolumn{2}{c|}{\textbf{Mixed}} & \multicolumn{2}{c|}{\textbf{Tensor}} \\
    \hline
    $S_0$ & $[-4.96,\,4.90]$ & $M_0$ & $[-2.08,\,2.22]$ & $T_0$ & $[-0.32,\,0.22]$ \\
    $S_1$ & $-$              & $M_1$ & $[-5.44,\,5.18]$ & $T_1$ & $[-0.16,\,0.16]$ \\
          &                  & $M_2$ & $[-23.7,\,22.8]$ & $T_2$ & $[-0.61,\,0.42]$ \\
          &                  & $M_3$ & $[-6.56,\,6.63]$ & $T_5$ & $[-0.15,\,0.17]$ \\
          &                  & $M_4$ & $[-3.59,\,3.63]$ & $T_6$ & $[-0.55,\,0.57]$ \\
          &                  & $M_5$ & $[-3.04,\,3.10]$ & $T_7$ & $[-0.52,\,0.56]$ \\
          &                  & $M_7$ & $[-7.90,\,7.65]$ & $T_8$ & $[-0.28,\,0.28]$ \\
          &                  &       &                  & $T_9$ & $[-0.53,\,0.53]$ \\
    \hline
  \end{tabular}
\end{table}

The fully hadronic channel pursued in this thesis carries the largest branching fraction of all VBS topologies, approximately $46\%$ for $WW$, $47\%$ for $WZ$ and $49\%$ for $ZZ$ in the four-quark final state~\cite{ParticleDataGroup:2024cfk}. Its statistical reach therefore complements rather than duplicates the leptonic channels, and including the present analysis in a combination of the kind of Ref.~\cite{ATLAS:VBStribosoncombination:2026} is the natural route to tightening the per-operator limits beyond any single channel. Such a combination requires consistent EFT-model assumptions across the inputs, in particular a common treatment of the linear interference term (Section~\ref{sec:disc_interference}) and a common unitarity prescription~\cite{PhysRevD.101.113003}.

\section{Impact of the Run-3 dataset}
\label{sec:outlook_run3}

The analysis is designed as a Run-2 plus Run-3 measurement. Run-3, now complete, was recorded between 2022 and 2026 at the higher centre-of-mass energy of $\sqrt{s}=13.6\,\mathrm{TeV}$ and yielded $332\,\mathrm{fb}^{-1}$ recorded in ATLAS, of $351\,\mathrm{fb}^{-1}$ delivered~\cite{ATLAS:lumiRun3}, more than twice the $140\,\mathrm{fb}^{-1}$ of Run-2 analysed here and at a slightly larger VBS cross section. Re-running the analysis chain on the Run-3 data is the immediate way to improve the constraints, with the larger gain expected for the statistics-limited scalar and mixed families and a smaller one for the tensor families, whose high-mass tail already approaches the systematics-limited regime (Section~\ref{sec:sens_tensor}). A quantitative projection of the combined reach is deliberately not attempted here: it requires the Run-3 calibration, dedicated $13.6\,\mathrm{TeV}$ signal samples and the full systematic model of Section~\ref{sec:disc_systematics}, none of which is available at this stage.
